% chapters/requirements/rf_assignment.tex
% RF-8: Reglas de corrección
% ============================================================================

\item \textbf{Reglas de corrección.} \label{rf:assignment}
El coordinador debe poder definir reglas de asignación (\texttt{AssignmentRule}) que
determinen qué profesores corrigen qué instancias o problemas. Las reglas se basan en
filtros combinables: grupo, modelo y problema. El sistema soporta que varios correctores
estén asignados al mismo problema de la misma instancia; en ese caso, todos pueden
calificar y la coordinación se resuelve de forma externa.

\begin{functional}

    % RF-8.1 ----------------------------------------------------------------
    \item \textbf{Creación de reglas de asignación.} \label{rf:assignment:create}
    El sistema debe permitir al coordinador crear reglas que asignen correctores a
    instancias o problemas.

    \textbf{Entrada:} petición POST al recurso de reglas del examen con los datos de la
    regla: lista de profesores correctores y filtros de alcance combinables: grupo (uno,
    varios o todos), modelo (uno, varios o todos) y problema (uno, varios o todos).

    \textbf{Proceso:} el sistema almacena la regla y calcula las asignaciones concretas que
    genera: qué corrector evalúa qué instancia y qué problemas. Por ejemplo, una regla
    <<el profesor~A corrige los problemas~1, 3 y~5 del modelo~A>> genera una asignación
    por cada instancia del modelo~A.

    \textbf{Salida:} respuesta con la regla creada, el número de asignaciones generadas y
    código \acs{http}~201.

    % RF-8.2 ----------------------------------------------------------------
    \item \textbf{Modificación de reglas de asignación.} \label{rf:assignment:update}
    El sistema debe permitir al coordinador modificar una regla de asignación existente.

    \textbf{Entrada:} petición PATCH al recurso de la regla con los datos a modificar.

    \textbf{Proceso:} el sistema recalcula las asignaciones derivadas de la regla
    modificada. Si la modificación elimina asignaciones sobre instancias que ya tienen
    calificaciones, se emite una advertencia pero se permite la modificación.

    \textbf{Salida:} respuesta con la regla actualizada, el nuevo número de asignaciones y,
    si procede, la advertencia sobre calificaciones existentes.

    % RF-8.3 ----------------------------------------------------------------
    \item \textbf{Eliminación de reglas de asignación.} \label{rf:assignment:delete}
    El sistema debe permitir al coordinador eliminar una regla de asignación.

    \textbf{Entrada:} petición DELETE al recurso de la regla.

    \textbf{Proceso:} el sistema elimina la regla y las asignaciones derivadas. Si hay
    calificaciones asociadas a asignaciones que se eliminan, estas calificaciones se
    conservan pero los correctores pierden el acceso a esas instancias.

    \textbf{Salida:} respuesta con código \acs{http}~204.

    % RF-8.4 ----------------------------------------------------------------
    \item \textbf{Aplicación de reglas y cálculo de asignaciones.}
    \label{rf:assignment:apply}
    El sistema debe aplicar todas las reglas activas del examen para generar las
    asignaciones concretas de corrector a instancias y problemas.

    \textbf{Entrada:} petición POST al recurso de aplicación de reglas del examen, o evento
    interno tras la creación de nuevas instancias (ingesta).

    \textbf{Proceso:} el sistema evalúa todas las reglas del examen contra el conjunto
    actual de instancias y genera o actualiza las asignaciones. Si una instancia nueva
    cumple los filtros de una regla existente, se le asigna automáticamente el corrector
    correspondiente.

    \textbf{Salida:} respuesta con el resumen de asignaciones generadas o actualizadas.

    % RF-8.5 ----------------------------------------------------------------
    \item \textbf{Verificación de cobertura.} \label{rf:assignment:coverage}
    El sistema debe verificar que las reglas definidas cubren todas las instancias y
    problemas del examen.

    \textbf{Entrada:} petición GET al recurso de verificación de cobertura del examen.

    \textbf{Proceso:} el sistema identifica las instancias y los problemas que no tienen al
    menos un corrector asignado. Agrupa los resultados por modelo y grupo para facilitar la
    identificación de huecos.

    \textbf{Salida:} respuesta indicando si la cobertura es completa o, en caso contrario,
    la lista detallada de instancias y problemas sin corrector asignado.

    % RF-8.6 ----------------------------------------------------------------
    \item \textbf{Listado de tareas del corrector.} \label{rf:assignment:tasks}
    El sistema debe permitir a un corrector consultar las instancias y problemas que tiene
    asignados para corregir.

    \textbf{Entrada:} petición GET al recurso de tareas del corrector, con filtro opcional
    de examen y estado de corrección.

    \textbf{Proceso:} el sistema recupera todas las asignaciones activas del corrector
    autenticado (derivadas de \texttt{AssignmentRule}), agrupadas por examen. Para cada
    asignación indica: instancia, problemas asignados, estado de corrección de cada
    problema (calificado o pendiente) y si la instancia tiene incidencias.

    \textbf{Salida:} respuesta con la lista de tareas del corrector, paginada y con los
    indicadores de estado.

\end{functional}
